% ---------------------------------------------------------------
% seccoes/conclusao.tex — Conclusão
% ---------------------------------------------------------------

\chapter{Conclusão}\label{sec:8}

Em resposta à questão central apresentada na Secção~\ref{sec:1.2}, o trabalho desenvolvido demonstra que é possível utilizar um motor de regras associado ao processamento de eventos complexos para apoiar a tomada de decisão de um veículo, em tempo real, a partir de dados provenientes de sensores com diferentes frequências de atualização. Para tal, foi fundamental estruturar o processo de decisão em três etapas distintas: perceção, geração de propostas e decisões.

A base de conhecimento, construída com base na regulamentação rodoviária, na literatura e no contributo do perito do domínio, integra um conjunto de regras distintas e contempla cenários predefinidos de condução. A combinação do raciocínio temporal e da lógica difusa permitiu abordar diferentes situações de circulação e incorporar distintas formas de raciocínio no processo de decisão. A distinção entre risco e prioridade, aliada à decisão em duas vertentes (velocidade e via), permite que a escolha das ações não dependa da ordem de execução das regras.

Neste sentido, o principal valor do presente trabalho diz respeito à abordagem na tomada de decisão transparente e rastreável, na qual os comportamentos podem ser especificados, testados, analisados e justificados. Embora a solução ainda apresente limitações próprias de um protótipo académico, constitui uma base para desenvolvimentos futuros, particularmente ao nível da fiabilidade dos sensores, do tratamento da incerteza e da interação entre veículos, aproximando progressivamente a abordagem das exigências de sistemas de condução autónoma em contextos reais.
